\documentclass[11pt,a4paper]{article}

% ===== PACKAGES =====
\usepackage[utf8]{inputenc}
\usepackage[T1]{fontenc}
\usepackage[margin=1in]{geometry}
\usepackage{hyperref}
\usepackage{listings}
\usepackage{xcolor}
\usepackage{booktabs}
\usepackage{array}
\usepackage{longtable}
\usepackage{fancyhdr}
\usepackage{titlesec}
\usepackage{tcolorbox}
\usepackage{fontawesome5}
\usepackage{amssymb}
\usepackage{amsmath}

% ===== COLORS =====
\definecolor{codegreen}{rgb}{0.25,0.49,0.48}
\definecolor{codegray}{rgb}{0.5,0.5,0.5}
\definecolor{codepurple}{rgb}{0.58,0,0.82}
\definecolor{codeblue}{rgb}{0.13,0.29,0.53}
\definecolor{codeorange}{rgb}{0.8,0.4,0.0}
\definecolor{backcolour}{rgb}{0.95,0.95,0.95}
\definecolor{warningbg}{rgb}{1.0,0.95,0.85}
\definecolor{warningborder}{rgb}{0.9,0.6,0.2}
\definecolor{tipbg}{rgb}{0.9,0.97,0.9}
\definecolor{tipborder}{rgb}{0.3,0.7,0.3}

% ===== CODE LISTING STYLE =====
\lstdefinelanguage{IEC61131}{
    keywords={VAR, END_VAR, VAR_INPUT, VAR_OUTPUT, VAR_IN_OUT, VAR_GLOBAL, 
              TYPE, END_TYPE, STRUCT, END_STRUCT, FUNCTION, END_FUNCTION,
              FUNCTION_BLOCK, END_FUNCTION_BLOCK, PROGRAM, END_PROGRAM,
              IF, THEN, ELSE, ELSIF, END_IF, CASE, OF, END_CASE,
              FOR, TO, BY, DO, END_FOR, WHILE, END_WHILE, REPEAT, UNTIL, END_REPEAT,
              TRUE, FALSE, AND, OR, NOT, XOR, MOD,
              ARRAY, OF, POINTER, REFERENCE, TO, AT},
    keywordstyle=\color{codeblue}\bfseries,
    keywords=[2]{BOOL, BYTE, WORD, DWORD, LWORD, SINT, USINT, INT, UINT, 
                 DINT, UDINT, LINT, ULINT, REAL, LREAL, STRING, WSTRING,
                 TIME, LTIME, DATE, TIME_OF_DAY, TOD, DATE_AND_TIME, DT},
    keywordstyle=[2]\color{codepurple}\bfseries,
    keywords=[3]{ADR, SIZEOF, REF, __ISVALIDREF},
    keywordstyle=[3]\color{codeorange}\bfseries,
    sensitive=false,
    morecomment=[l]{//},
    morecomment=[s]{(*}{*)},
    commentstyle=\color{codegreen}\itshape,
    stringstyle=\color{codeorange},
    morestring=[b]',
    morestring=[b]",
}

\lstdefinestyle{iecstyle}{
    language=IEC61131,
    backgroundcolor=\color{backcolour},
    basicstyle=\ttfamily\small,
    breakatwhitespace=false,
    breaklines=true,
    captionpos=b,
    keepspaces=true,
    numbers=left,
    numbersep=8pt,
    numberstyle=\tiny\color{codegray},
    showspaces=false,
    showstringspaces=false,
    showtabs=false,
    tabsize=4,
    frame=single,
    framerule=0.5pt,
    rulecolor=\color{codegray},
    xleftmargin=15pt,
    framexleftmargin=15pt,
}

\lstset{style=iecstyle}

% ===== TCOLORBOX STYLES =====
\tcbuselibrary{skins,breakable}

\newtcolorbox{warningbox}{
    colback=warningbg,
    colframe=warningborder,
    boxrule=1pt,
    left=10pt,
    right=10pt,
    top=8pt,
    bottom=8pt,
    breakable,
    before upper={\faExclamationTriangle\ \textbf{Warning:} }
}

\newtcolorbox{tipbox}{
    colback=tipbg,
    colframe=tipborder,
    boxrule=1pt,
    left=10pt,
    right=10pt,
    top=8pt,
    bottom=8pt,
    breakable,
    before upper={\faLightbulb\ \textbf{Tip:} }
}

% ===== HEADER/FOOTER =====
\pagestyle{fancy}
\fancyhf{}
\fancyhead[L]{\leftmark}
\fancyhead[R]{TwinCAT 3 Lecture Notes}
\fancyfoot[C]{\thepage}
\renewcommand{\headrulewidth}{0.4pt}
\renewcommand{\footrulewidth}{0.4pt}

% ===== HYPERREF SETUP =====
\hypersetup{
    colorlinks=true,
    linkcolor=codeblue,
    filecolor=magenta,
    urlcolor=cyan,
    pdftitle={TwinCAT 3 - Data Types, Pointers, and Arrays},
    pdfauthor={},
    bookmarks=true,
}

% ===== TITLE FORMATTING =====
\titleformat{\section}
    {\normalfont\Large\bfseries}{\thesection}{1em}{}[{\titlerule[0.8pt]}]
\titleformat{\subsection}
    {\normalfont\large\bfseries}{\thesubsection}{1em}{}
\titleformat{\subsubsection}
    {\normalfont\normalsize\bfseries}{\thesubsubsection}{1em}{}

% ===== DOCUMENT INFO =====
\title{
    \vspace{-1cm}
    \textbf{Lecture Notes: TwinCAT 3}\\
    \Large Data Types, Pointers, and Arrays
}
\author{}
\date{\today}

% ===== BEGIN DOCUMENT =====
\begin{document}

\maketitle

\tableofcontents
\newpage

% ============================================================
\section{Introduction to Variables and Static Typing}
% ============================================================

\subsection{Variables as Containers}

Variables label and store data in memory so programs are human-readable and data can be manipulated.

\subsection{Static vs. Dynamic Typing}

Unlike JavaScript or Python (dynamic), \textbf{IEC 61131-3 is statically typed}:

\begin{table}[h]
\centering
\begin{tabular}{@{}lll@{}}
\toprule
\textbf{Aspect} & \textbf{Static Typing (IEC 61131-3)} & \textbf{Dynamic Typing (JS/Python)} \\
\midrule
Type Checking & Compile time & Runtime \\
Type Declaration & Required before use & Not required \\
Type Changes & Not allowed during execution & Allowed \\
\bottomrule
\end{tabular}
\caption{Static vs. Dynamic Typing Comparison}
\end{table}

\subsection{Declaration Block}

In TwinCAT 3, variables are declared between the keywords \texttt{VAR} and \texttt{END\_VAR}.

\subsubsection*{Example: Basic Declaration}

\begin{lstlisting}
VAR
    fCabinetTemperature : REAL;  // Declares a 4-byte floating point variable
END_VAR
\end{lstlisting}

% ============================================================
\section{Standard Data Types}
% ============================================================

\subsection{Integer and Bit-Oriented Types}

\subsubsection{Integer Sizes}

\begin{table}[h]
\centering
\begin{tabular}{@{}llllll@{}}
\toprule
\textbf{Type} & \textbf{Size} & \textbf{Signed Range} & \textbf{Unsigned Type} & \textbf{Unsigned Range} \\
\midrule
SINT/USINT & 1 byte & -128 to 127 & USINT & 0 to 255 \\
INT/UINT & 2 bytes & -32,768 to 32,767 & UINT & 0 to 65,535 \\
DINT/UDINT & 4 bytes & $-2^{31}$ to $2^{31}-1$ & UDINT & 0 to $2^{32}-1$ \\
LINT/ULINT & 8 bytes & $-2^{63}$ to $2^{63}-1$ & ULINT & 0 to $2^{64}-1$ \\
\bottomrule
\end{tabular}
\caption{Integer Data Types}
\end{table}

\subsubsection{Base Notation}

You can define integers using different bases:

\begin{table}[h]
\centering
\begin{tabular}{@{}lll@{}}
\toprule
\textbf{Base} & \textbf{Prefix} & \textbf{Example} \\
\midrule
Binary & \texttt{2\#} & \texttt{2\#1011} \\
Octal & \texttt{8\#} & \texttt{8\#17} \\
Hexadecimal & \texttt{16\#} & \texttt{16\#FF} \\
\bottomrule
\end{tabular}
\caption{Base Notation Prefixes}
\end{table}

\subsubsection{Bit Accessing}

Access individual bits using dot notation: \texttt{nVariable.0} (first bit)

\subsubsection*{Example: Integer Bases and Bit Access}

\begin{lstlisting}
VAR
    byBinaryValue  : BYTE := 2#11001010;   // Binary notation (decimal 202)
    wHexValue      : WORD := 16#FF0A;      // Hexadecimal notation (decimal 65290)
    nStatusFlags   : BYTE;
    bBitThree      : BOOL;
END_VAR

// Accessing the 4th bit (index 3) of nStatusFlags
bBitThree := nStatusFlags.3;
\end{lstlisting}

\subsection{Floating Point and Strings}

\begin{table}[h]
\centering
\begin{tabular}{@{}llll@{}}
\toprule
\textbf{Type} & \textbf{Size} & \textbf{Equivalent} & \textbf{Description} \\
\midrule
\texttt{REAL} & 4 bytes & Float & Single precision \\
\texttt{LREAL} & 8 bytes & Double & Double precision \\
\texttt{STRING} & 81 bytes (default) & --- & ASCII-coded, 80 chars + null terminator \\
\texttt{WSTRING} & Variable & --- & UNICODE/UTF-16, 2--4 bytes per character \\
\bottomrule
\end{tabular}
\caption{Floating Point and String Types}
\end{table}

\subsection{Time and Date}

TwinCAT provides specific types to handle time in human-readable formats:

\begin{table}[h]
\centering
\begin{tabular}{@{}lllll@{}}
\toprule
\textbf{Type} & \textbf{Resolution} & \textbf{Prefix} & \textbf{Example} \\
\midrule
\texttt{TIME} & 1ms & \texttt{T\#} & \texttt{T\#2s}, \texttt{T\#100ms} \\
\texttt{LTIME} & 1ns & \texttt{LTIME\#} & \texttt{LTIME\#500us} \\
\texttt{DATE} & 1 day & \texttt{D\#} & \texttt{D\#2024-01-15} \\
\texttt{TIME\_OF\_DAY} & 1ms & \texttt{TOD\#} & \texttt{TOD\#14:30:00} \\
\texttt{DATE\_AND\_TIME} & 1ms & \texttt{DT\#} & \texttt{DT\#2024-01-15-14:30:00} \\
\bottomrule
\end{tabular}
\caption{Time and Date Types}
\end{table}

% ============================================================
\section{User-Defined Types: Enumerations (ENUMs)}
% ============================================================

\subsection{Definition}

A set of named integral constants used to make code self-explanatory.

\subsection{Key Points}

\begin{itemize}
    \item \textbf{Scope:} ENUMs are \textbf{globally accessible} in TwinCAT
    \item \textbf{Risk:} Can lead to ``namespace pollution'' if multiple libraries use the same ENUM names
    \item \textbf{Creation:} Created as a \textbf{DUT (Data Unit Type)} in the project tree
\end{itemize}

\subsubsection*{Example: DUT Definition}

\begin{lstlisting}
// DUT Definition (created as a separate DUT file in the project tree)
{attribute 'qualified_only'}
TYPE E_TrafficLight :
(
    Red    := 0,
    Yellow := 1,
    Green  := 2
);
END_TYPE
\end{lstlisting}

\subsubsection*{Example: Usage in Program}

\begin{lstlisting}
// Usage in a Program or Function Block
VAR
    eCurrentLight : E_TrafficLight;
END_VAR

// Assigning using the dot identifier (required with 'qualified_only' attribute)
eCurrentLight := E_TrafficLight.Green;
\end{lstlisting}

\begin{tipbox}
Use the \texttt{\{attribute 'qualified\_only'\}} pragma to require the full enum name (e.g., \texttt{E\_TrafficLight.Green}) and avoid namespace conflicts.
\end{tipbox}

% ============================================================
\section{Memory Operations and SIZEOF}
% ============================================================

\subsection{Memory Allocation}

\begin{itemize}
    \item Memory is generally allocated \textbf{upfront} at compile time
    \item Dynamic allocation is typically \textbf{discouraged} in PLC programming
\end{itemize}

\subsection{SIZEOF Operator}

Used to determine the memory size (in bytes) of a variable or type.

\begin{table}[h]
\centering
\begin{tabular}{@{}ll@{}}
\toprule
\textbf{Target} & \textbf{Returns} \\
\midrule
\texttt{STRING} (default) & 81 bytes \\
Reference & Size of the \textbf{referenced object} \\
Pointer & Size of the \textbf{pointer itself} (e.g., 8 bytes on 64-bit) \\
\bottomrule
\end{tabular}
\caption{SIZEOF Operator Behavior}
\end{table}

% ============================================================
\section{Pointers vs. References}
% ============================================================

\subsection{Pointers}

Pointers represent a memory address.

\subsubsection{Operators}

\begin{table}[h]
\centering
\begin{tabular}{@{}lll@{}}
\toprule
\textbf{Operator} & \textbf{Purpose} & \textbf{Example} \\
\midrule
\texttt{ADR()} & Get address of variable & \texttt{pVar := ADR(nValue)} \\
\texttt{\^{}} & Dereference (access value) & \texttt{pVar\^{} := 100} \\
\bottomrule
\end{tabular}
\caption{Pointer Operators}
\end{table}

\subsubsection{Risks}

\begin{warningbox}
Pointers are \textbf{not type-safe}. You can accidentally point a STRING pointer at an INT variable, resulting in ``garbage'' data.
\end{warningbox}

\subsubsection*{Example: Pointer Usage}

\begin{lstlisting}
VAR
    nMyValue    : INT := 100;
    pnMyPointer : POINTER TO INT;
END_VAR

// Assign the address of nMyValue to the pointer
pnMyPointer := ADR(nMyValue);

// Dereference the pointer to update the value
pnMyPointer^ := 250;  // nMyValue is now 250
\end{lstlisting}

\subsection{References}

References act as an \textbf{alias} for an object.

\subsubsection{Operators}

\begin{table}[h]
\centering
\begin{tabular}{@{}ll@{}}
\toprule
\textbf{Operator} & \textbf{Purpose} \\
\midrule
\texttt{REFERENCE TO} & Declare a reference \\
\texttt{REF=} & Assign a reference \\
\bottomrule
\end{tabular}
\caption{Reference Operators}
\end{table}

\subsubsection{Advantages}

\begin{table}[h]
\centering
\begin{tabular}{@{}ll@{}}
\toprule
\textbf{Feature} & \textbf{Benefit} \\
\midrule
\textbf{Type Safety} & Compiler prevents assigning to wrong data type \\
\textbf{Syntax} & No special dereferencing (\texttt{\^{}}) needed \\
\bottomrule
\end{tabular}
\caption{Reference Advantages}
\end{table}

\subsubsection{Safety Check}

Use \texttt{\_\_ISVALIDREF()} to check if a reference points to a valid memory location before use.

\begin{lstlisting}
IF __ISVALIDREF(refMyVariable) THEN
    // Safe to use the reference
END_IF
\end{lstlisting}

% ============================================================
\section{Arrays}
% ============================================================

\subsection{Definition}

A series of elements of the \textbf{same type} in \textbf{contiguous memory}.

\subsection{Key Features}

\begin{table}[h]
\centering
\begin{tabular}{@{}ll@{}}
\toprule
\textbf{Feature} & \textbf{Description} \\
\midrule
\textbf{Flexible Indices} & Can start/end at any integer (e.g., \texttt{ARRAY [1..5]}, \texttt{ARRAY [-1..1]}) \\
\textbf{Multi-Dimensional} & Support for 2D, 3D, etc. arrays \\
\bottomrule
\end{tabular}
\caption{Array Features}
\end{table}

\subsubsection*{Example: Array Declaration and Access}

\begin{lstlisting}
VAR
    // 1D Array: 5 integers with indices 0 through 4
    anValues : ARRAY [0..4] OF INT := [10, 20, 30, 40, 50];
    
    // 2D Array - Alternative #1: Explicit row and column indices
    anMatrix1 : ARRAY [0..2, 0..3] OF INT;  // 3 rows, 4 columns
    
    // 2D Array - Alternative #2: Array of arrays
    anMatrix2 : ARRAY [0..2] OF ARRAY [0..3] OF INT;  // 3 rows, 4 columns
END_VAR

// Accessing elements
anValues[2] := 35;           // Set the 3rd element to 35
anMatrix1[1, 2] := 99;       // Set row 1, column 2 to 99
anMatrix2[1][2] := 99;       // Equivalent access for array of arrays
\end{lstlisting}

\subsection{2D Array Declaration Comparison}

\begin{table}[h]
\centering
\begin{tabular}{@{}lll@{}}
\toprule
\textbf{Style} & \textbf{Syntax} & \textbf{Access Pattern} \\
\midrule
Alternative \#1 & \texttt{ARRAY [0..2, 0..3] OF INT} & \texttt{arr[row, col]} \\
Alternative \#2 & \texttt{ARRAY [0..2] OF ARRAY [0..3] OF INT} & \texttt{arr[row][col]} \\
\bottomrule
\end{tabular}
\caption{2D Array Declaration Styles}
\end{table}

% ============================================================
\section{Type Conversion}
% ============================================================

\subsection{Implicit vs. Explicit Conversion}

\begin{table}[h]
\centering
\begin{tabular}{@{}lllll@{}}
\toprule
\textbf{Direction} & \textbf{Type} & \textbf{Example} & \textbf{Operator Required} \\
\midrule
Smaller $\rightarrow$ Larger & Implicit & USINT $\rightarrow$ UINT & No \\
Larger $\rightarrow$ Smaller & Explicit & UINT $\rightarrow$ USINT & Yes (\texttt{UINT\_TO\_USINT}) \\
\bottomrule
\end{tabular}
\caption{Type Conversion Rules}
\end{table}

\subsection{Data Loss Warning}

\begin{warningbox}
Converting a large value to a smaller type truncates the high bytes.
\end{warningbox}

\textbf{Example:}

\begin{verbatim}
Decimal 259 = Binary 00000001 00000011 (2 bytes)
                            |
                            v
Truncated to USINT = Binary 00000011 = Decimal 3
\end{verbatim}

\subsection{Common Conversion Operators}

\begin{table}[h]
\centering
\begin{tabular}{@{}lll@{}}
\toprule
\textbf{From} & \textbf{To} & \textbf{Operator} \\
\midrule
INT & REAL & \texttt{INT\_TO\_REAL()} \\
REAL & INT & \texttt{REAL\_TO\_INT()} \\
UINT & USINT & \texttt{UINT\_TO\_USINT()} \\
STRING & INT & \texttt{STRING\_TO\_INT()} \\
\bottomrule
\end{tabular}
\caption{Common Type Conversion Operators}
\end{table}

% ============================================================
\newpage
\section*{Quick Reference Card}
\addcontentsline{toc}{section}{Quick Reference Card}
% ============================================================

\subsection*{Declaration Syntax}

\begin{lstlisting}
VAR
    <name> : <TYPE>;                           // Basic
    <name> : <TYPE> := <initial_value>;        // With initialization
    <name> : ARRAY [<start>..<end>] OF <TYPE>; // Array
    <name> : POINTER TO <TYPE>;                // Pointer
    <name> : REFERENCE TO <TYPE>;              // Reference
END_VAR
\end{lstlisting}

\subsection*{Common Prefixes (Hungarian Notation)}

\begin{table}[h]
\centering
\begin{tabular}{@{}lll@{}}
\toprule
\textbf{Prefix} & \textbf{Type} & \textbf{Example} \\
\midrule
\texttt{b} & BOOL & \texttt{bIsRunning} \\
\texttt{n} & INT/UINT/etc. & \texttt{nCounter} \\
\texttt{f} & REAL/LREAL & \texttt{fTemperature} \\
\texttt{s} & STRING & \texttt{sMessage} \\
\texttt{a} & ARRAY & \texttt{anValues} \\
\texttt{p} & POINTER & \texttt{pnData} \\
\texttt{ref} & REFERENCE & \texttt{refMotor} \\
\texttt{e} & ENUM & \texttt{eState} \\
\bottomrule
\end{tabular}
\caption{Hungarian Notation Prefixes}
\end{table}

\vfill

\begin{center}
\textit{Last Updated: \today}
\end{center}

\end{document}